% Generated draft: requires mathematical and source review.
\documentclass[11pt]{article}
\usepackage[margin=25mm]{geometry}
\usepackage{fontspec}
\setmainfont{Latin Modern Roman}
\usepackage{amsmath,amssymb,mathtools,unicode-math}
\setmathfont{Latin Modern Math}
\usepackage{enumitem,needspace,xurl,hyperref}
\setlength{\parindent}{0pt}
\setlength{\parskip}{5pt}
\setlength{\emergencystretch}{3em}
\newcommand{\N}{\mathbb{N}}
\newcommand{\Z}{\mathbb{Z}}
\newcommand{\Q}{\mathbb{Q}}
\newcommand{\R}{\mathbb{R}}
\newcommand{\C}{\mathbb{C}}
\newcommand{\F}{\mathbb{F}}
\newcommand{\A}{\mathbb{A}}
\newcommand{\PP}{\mathbb{P}}
\newcommand{\E}{\mathbb{E}}
\newcommand{\eps}{\varepsilon}
\newcommand{\dd}{\,\mathrm{d}}
\begin{document}
\section*{UnsolvedMath ID: 30005844}
\textbf{Source problem code:} catalogue entry
\subsection{Definitions and mathematical statement}
Use the standard convention that matrix entries are nonnegative integers. For positive integers \(n,m\), a square matrix is called essential if it has no zero row and no zero column. Matrices \(A\in M_n(\Z_{\geq 0})\) and \(B\in M_m(\Z_{\geq 0})\) are shift equivalent over the nonnegative integers if there exist \(\ell\geq 1\), \(R\in M_{n\times m}(\Z_{\geq 0})\), and \(S\in M_{m\times n}(\Z_{\geq 0})\) such that
\[
 AR=RB,\qquad SA=BS,\qquad A^\ell=RS,\qquad B^\ell=SR.
\]
For a common coefficient field \(K\), let \(E_A\) be the finite directed graph with vertices \(1,\ldots,n\) and exactly \(A_{ij}\) edges from \(i\) to \(j\). Write \(L_K(E_A)\) for its Leavitt path algebra. Explicitly, it is the \(K\)-algebra generated by vertices \(v\), edges \(e\), and formal ghost edges \(e^*\), subject to
\[
 uv=\delta_{u,v}u,\qquad s(e)e=e=e r(e),\qquad r(e)e^*=e^*=e^*s(e),
\]
\[
 e^*f=\delta_{e,f}r(e),\qquad
 v=\sum_{s(e)=v}ee^*.
\]
The last relation is imposed at each vertex; essentiality ensures that each such sum is finite and nonempty. Its standard \(\Z\)-grading is specified by \(\deg(v)=0\), \(\deg(e)=1\), and \(\deg(e^*)=-1\). Define \(E_B\) and \(L_K(E_B)\) in the same way.

For a graded algebra \(D\), let \(\operatorname{GrMod}(D)\) denote the category of unital \(\Z\)-graded right \(D\)-modules with degree-preserving homomorphisms. Adopt the shift convention \(M(t)_d=M_{d+t}\). In the standard interpretation used here, graded Morita equivalence means an equivalence of these graded module categories compatible, up to coherent natural isomorphisms, with every integer grading shift.

\textbf{Conjecture 2 (Hazrat), as supplied.} For all essential square matrices \(A\) and \(B\) over the nonnegative integers, is
\[
 A\text{ and }B\text{ shift equivalent}
 \quad\Longleftrightarrow\quad
 L_K(E_A)\text{ and }L_K(E_B)\text{ graded Morita equivalent}?
\]
Here \(L_K(E_A)\) and \(L_K(E_B)\) make explicit the usual interpretation of the source notation \(L_A,L_B\). The supplied record does not specify the coefficient field or its intended quantification. Consequently, this is a conditional formulation for a common field \(K\); whether the intended conjecture fixes a particular field or quantifies over all fields remains unresolved. The definitions of essentiality, graph orientation, and graded Morita equivalence above state conventional interpretations that the supplied excerpt does not itself define.
\subsection{Short English statement}
The conjecture attributed to Hazrat asks whether shift equivalence of essential nonnegative integer matrices is exactly equivalent to graded Morita equivalence of their associated Leavitt path algebras. The coefficient field and its intended quantification are missing from the supplied statement. The supplied dated curation labels the problem partially solved; this draft does not independently verify that assessment.
\subsection{Sources}
\begin{itemize}
\item Local UnsolvedMath registry, permanent ID \(30005844\), title \emph{Hazrat Conjecture for Leavitt Path Algebras}, exact source problem code \texttt{catalogue entry}. The supplied excerpt attributes Conjecture 2 to Hazrat; the registry's proposer field is empty.
\item \emph{Combinatorial \(*\)-algebras} (2024), Oberwolfach Reports, pp. 731--788, DOI 10.4171/owr/2024/12. Supplied workshop-report provenance: \url{https://doi.org/10.4171/owr/2024/12}. The record does not identify the particular contribution, author, or page containing the conjecture.
\item \emph{Recasting the Hazrat conjecture}. Supplied literature reference: \url{https://doi.org/10.1007/s10468-024-10266-w}.
\item \emph{Unital aligned shift equivalence and graded classification}. Supplied literature reference: \url{https://doi.org/10.1007/s00209-025-03748-0}.
\item Literature assessment dated exactly \texttt{2026-08-22}: supplied dated curation, not independently verified. It labels the problem partially solved and states that the Hazrat conjecture remains unresolved in full generality, while aligned or unital shift equivalences and additional compatibility hypotheses imply graded isomorphism or graded Morita equivalence in important classes. The record explicitly requires independent verification of this time-sensitive status and the cited claims.
\end{itemize}
\end{document}
